\documentclass[11pt,a4paper]{article}
\usepackage{lmodern}
\usepackage[utf8]{inputenc}
\usepackage[T1]{fontenc}
\usepackage{amsmath}
\usepackage{amsfonts}
\usepackage{amssymb}
\usepackage{graphicx}
\usepackage[french]{babel}
\usepackage[notitlepage,fancysections,pagenumber]{sauveterre}

\usepackage{enumitem}

\usepackage{indentfirst}
\usepackage{lipsum}

\usepackage{cases}
\usepackage{multicol}

%formatting titles
\usepackage{titlesec}
\titleformat*{\section}{\centering\Large\bfseries}
\titleformat*{\subsection}{\centering\large\bfseries}

%formatting date
\usepackage[ddmmyyyy]{datetime}
\renewcommand{\dateseparator}{/}

%\usepackage{geometry}
%\geometry{
% a4paper,
% total={170mm,257mm},
% left=20mm,
% top=20mm,
% }

\usepackage{pgf,tikz}
\usepackage{mathrsfs}
\usetikzlibrary{arrows}

\usepackage[export]{adjustbox}

\renewcommand{\notitlepageupperline}{
	\noindent%
	\begin{minipage}{\textwidth}
		\centering
		\begin{minipage}{0.33\textwidth}
			\raisebox{-0.5\height}
			{\color{bleu-sauveterre}%
				\includegraphics[height=0.8\hauteurlogovert]
				{\polylogovert}
				}
		\end{minipage}%
		\begin{minipage}{0.33\textwidth}
			\color{bleu-sauveterre}%
			\centering\polydatesave
			\polyauthorsave
			\hspace{3em}
		\end{minipage}%
		\begin{minipage}{0.33\textwidth}
			\color{bleu-sauveterre}%
			\flushright\raisebox{-0.5\height}
			{\includegraphics[height=0.8\hauteurlogovert]{logo-federal}
			}
		\end{minipage}
	\end{minipage}
}

%\usepackage{fancyhdr}
%\usepackage{lastpage}
%\pagestyle{fancy}
%\fancyhf{}
%%headers and footers
%\lhead{Centre Équestre de Sauveterre}
%\chead{Règlement intérieur}
%\rhead{Version du \today}
%\rfoot{Page \thepage \ sur \pageref{LastPage}}

%\fancypagestyle{firstpage}{%
%  \fancyhf{}
%  \renewcommand\headrulewidth{0pt}
%  \fancyfoot[R]{Page \thepage \ sur \pageref{LastPage}}
%}
%
%\title{\vspace{-50pt}
%\includegraphics[width=0.4\textwidth , left]{logo2-v2.eps}~
%\\[8pt]
%{\huge \textsc{Règlement intérieur}}
%}

\title{Mesures relatives à l'accès aux écuries en confinement}
\subtitle{Annexe temporaire III$^T$ du \textit{Règlement Intérieur} en date du \today}
\author{}
\date{}

%enumerate environment formatting
\renewcommand{\labelenumi}{\textbf{\arabic{enumi}.}}
\renewcommand{\labelenumii}{\textbf{\arabic{enumi}.\arabic{enumii}.}}
\renewcommand{\labelenumiii}{\textbf{\arabic{enumi}.\arabic{enumii}.\arabic{enumiii}.}}

%change sec and subsec numbering
\renewcommand{\thesection}{Section \Roman{section}.}
\renewcommand{\thesubsection}{Article \arabic{subsection} :}

\renewcommand{\labelenumi}{\arabic{enumi}.}


\renewcommand*\contentsname{Plan}

\renewcommand{\emph}[1]{\textit{\textbf{#1}}}


\usepackage{framed}

%\setlength{\parindent}{0pt}%remove indentation

\begin{document}

\maketitle

%\thispagestyle{firstpage}
%\vspace{-30pt}

Le présent document constitue une annexe temporaire au \textit{Règlement Intérieur}. Il entre en vigueur à la date de son affichage, révoque tout règlement de même référence précédemment établi, et sera abrogé dès lors qu'il sera mis à jour, ou ne sera plus affiché, ni à l'accueil ni sur le site internet du Centre Équestre.

\subsection*{Préambule}
	Suite aux mesures gouvernementales ayant pour objectif de limiter la propagation du COVID-19, et au couvre-feu avancé à 18h00 à compter du 16 janvier, l'accès aux écuries sera restreint au strict nécessaire au-delà de 18h00. Le présent règlement précise les conditions d'accès aux écuries et les mesures adoptées pour éviter au mieux la propagation du virus.
\subsection*{Engagement du Centre Équestre}
	Le Centre Équestre s'engage, par la présente, à maintenir les points suivants :
	\begin{itemize}
		\item mise à disposition de gel hydro-alcoolique à plusieurs emplacements stratégiques du club (accueil, devant le manège et à l'entrée des écuries) ;
		\item les visites des propriétaires seront prévues à l'avance pour limiter le nombre de personnes présentes dans les écuries ;
		\item les aires de travail, resteront disponibles selon les mêmes critères qu'avant ;
		\item les horaires de cours seront modifiés afin de permettre aux cavaliers de tous niveaux de continuer à monter au maximum ;
		\item la direction du Centre Équestre pourra fournir aux propriétaires d'équidés une attestation sur l'honneur permettant à ces derniers de venir s'en occuper, ainsi qu'une copie du contrat de pension si nécessaire ;
		\item l'équipe du Centre Équestre reste disponible en cas de besoin ou de problème avec les équidés, leur santé ou leur travail, le tout en respectant scrupuleusement les gestes barrière.
	\end{itemize}
	\subsection*{Engagements des propriétaires et cavaliers}
		En retour, les propriétaires et cavaliers s'engagent à respecter les points suivants :
		\begin{itemize}
			\item prévenir à l'avance de leur venue et du travail prévu, afin de gérer au mieux les flux et de limiter le nombre de personnes sur place ;
			\item respecter les consignes sanitaires telles que le port du masque dans l'enceinte du Centre Équestre ou encore la désinfection des mains par exemple ;
			\item venir après 18h00 \textbf{uniquement pour faire les soins de leur équidé}, après accord préalable de la direction ;
			\item ne pas rester plus longtemps que nécessaire sur place pour permettre à chacun de venir s'occuper de son équidé ;
			\item avertir l'équipe du Centre Équestre de tout problème rencontré ;
			\item ne pas se servir de l'attestation de pension ou du contrat de pension pour une autre raison que la venue au club pour le travail de son équidé.
		\end{itemize}

		De plus, les venues pendant le couvre-feu étant limitées aux soins, il est formellement interdit à tout cavalier de monter au-delà de 18h00. Toute personne allant à l'encontre de ce point est passible d'une exclusion à effet immédiat comme prévue par l'article 8 \textit{Règlement Intérieur}. Le Centre Équestre déclinant toute responsabilité quant aux propriétaires abusant des dérogations leur étant accordées et n'étant en aucun cas responsable en cas de verbalisation par les forces de l'ordre.


\vfill
\begin{flushright}
	Version du \today
\end{flushright}




\end{document}
